\chapter{The physics stack}
\label{ch:physics}

The tricycle runs on \textbf{Newton}, NVIDIA's GPU physics engine built on
Warp. Isaac Lab talks to Newton through a \emph{physics manager}, and the
tricycle uses two of Newton's solvers at once: a rigid-body solver for the car
and an implicit Material Point Method (MPM) solver for the soil. This chapter
explains the pieces that \code{tricycle_env_cfg.py} assembles, so that the
config in Chapter~\ref{ch:env_cfg} reads as a list of decisions rather than
magic numbers.

\section{Layers}

\begin{figure}[H]
\centering
\begin{tikzpicture}[
  layer/.style={draw,rounded corners=2pt,text width=13.2cm,minimum height=8mm,font=\footnotesize,align=center,fill=white},
  arr/.style={-{Stealth[length=1.8mm]},black!70}
]
\node[layer,fill=isaacblue!6,draw=isaacblue] (env) {\textbf{TricycleEnv / DirectRLEnv} \quad your methods + step bookkeeping};
\node[layer,below=3mm of env,fill=isaacblue!6,draw=isaacblue] (scene) {\textbf{InteractiveScene, Articulation, RigidObject, MPMObject} \quad tensors: \code{car.data.*}, \code{soil.data.*}};
\node[layer,below=3mm of scene,fill=isaacblue!6,draw=isaacblue] (sim) {\textbf{SimulationContext} \quad owns the stage and the physics manager};
\node[layer,below=3mm of sim,fill=newtonpurple!6,draw=newtonpurple] (mgr) {\textbf{NewtonCouplerManager} (\srcpath{isaaclab_contrib.coupling}) \quad builds one Newton Model, splits it into entries};
\node[layer,below=3mm of mgr,fill=newtonpurple!6,draw=newtonpurple] (solver) {\textbf{Newton SolverCoupledProxy} \quad entry ``rover'': SolverMuJoCo (MJWarp) \quad entry ``soil'': SolverImplicitMPM};
\node[layer,below=3mm of solver,fill=newtonpurple!6,draw=newtonpurple] (warp) {\textbf{Warp kernels on the GPU} \quad CUDA graph captured once, replayed every step};
\draw[arr] (env) -- (scene); \draw[arr] (scene) -- (sim); \draw[arr] (sim) -- (mgr); \draw[arr] (mgr) -- (solver); \draw[arr] (solver) -- (warp);
\end{tikzpicture}
\caption{From environment code to GPU kernels. Blue layers are Isaac Lab; purple layers are Newton.}
\end{figure}

\begin{isaacbox}{physics managers}
Isaac Lab develop can drive several physics backends (PhysX, Newton,
OmniPhysX). \srcpath{SimulationCfg.physics} selects one by holding a
\code{PhysicsCfg} subclass; \code{NewtonCfg}
(\srcpath{source/isaaclab_newton/isaaclab_newton/physics/newton_manager_cfg.py})
selects Newton. Inside \code{NewtonCfg}, \code{solver_cfg} selects
\emph{which} Newton solver manager runs the model: \code{MJWarpSolverCfg},
\code{MPMSolverCfg}, \code{XPBDSolverCfg}, \ldots, or, as here, a
\code{CouplerProxyCfg} from \srcpath{isaaclab_contrib.coupling} that nests two of
them. \srcpath{NewtonCfg.__post_init__} copies the manager class from
\srcpath{solver_cfg.class_type} so you never set it yourself.
\end{isaacbox}

\section{The rigid-body side: MuJoCo-Warp}

The car is an \emph{articulation}: a tree of rigid bodies (chassis, fork,
three wheels) connected by joints (one free joint for the chassis, one steering
hinge, three wheel hinges). Newton simulates it with \textbf{MJWarp}, a
GPU port of MuJoCo's dynamics. Isaac Lab wraps it as \code{MJWarpSolverCfg}
(\srcpath{physics/mjwarp_manager_cfg.py}). The four options the config sets:

\begin{description}[leftmargin=1.6em,font=\ttfamily]
  \item[use\_mujoco\_contacts=False] MuJoCo normally detects and resolves its
    own contacts. Setting this to \code{False} hands collision detection to
    Newton's own \emph{collision pipeline}, which is required whenever a
    Newton \code{collision_cfg} is also given (the two are mutually
    exclusive) and is what lets the coupler share contacts between solvers.
  \item[integrator="implicitfast"] MuJoCo's implicit-in-velocity integrator.
    It damps stiff joint drives and contacts far better than plain Euler at
    the same step size, which matters with a 10~ms step and velocity-driven
    wheels.
  \item[njmax=128] maximum number of constraint rows per env (joint limits,
    contacts, friction cones). Too small and contacts are dropped with a
    warning; too large only costs memory.
  \item[nconmax=128] maximum number of contact points per env.
\end{description}

\begin{isaacbox}{implicit actuators and \srcpath{use_newton_actuators}}
Our joints use \code{ImplicitActuatorCfg}. ``Implicit'' means the drive is a
PD controller \emph{inside the physics solver}: you give a position target
(steering) or a velocity target (rear wheels) plus \code{stiffness} and
\code{damping}, and the solver computes the torque as part of its implicit
step. \code{effort_limit_sim} and \code{velocity_limit_sim} are written into
the simulator as hard limits. \code{sim.use_newton_actuators = True} asks
Newton to author and execute these drives natively rather than through the
host-side adapter used for PhysX.
\end{isaacbox}

\section{The soil side: implicit MPM}
\label{sec:mpm}

\begin{newtonbox}{what MPM is}
The Material Point Method represents a continuum (soil, snow, mud) as a
cloud of \textbf{particles} that carry mass, velocity and deformation
history, plus a background \textbf{grid} on which forces are solved each
step. Every step: particle quantities are transferred to the grid
(``P2G''), the momentum equations with the material's constitutive law are
solved on the grid, and grid velocities are transferred back to move the
particles (``G2P''). Rigid bodies enter as \emph{colliders}: signed-distance
boundaries the grid velocities must respect. Newton's solver is
\emph{implicit}: it solves the yield condition as an optimisation problem
each step, which allows large steps (our 10~ms) without exploding.
\end{newtonbox}

Isaac Lab exposes it as \code{MPMSolverCfg}
(\srcpath{physics/mpm_manager_cfg.py}). The fields the tricycle sets, grouped:

\paragraph{Grid.}
\code{voxel_size=0.05} is the grid spacing (5~cm). \code{grid_type="sparse"}
allocates grid cells only where particles are. \code{grid_padding=0} adds no
empty margin. \srcpath{separate_worlds=True} gives every env its own grid so
envs cannot touch each other's soil and particle arrays are laid out
contiguously per env.

\paragraph{Sparse-grid capacities.}
\srcpath{max_active_cell_count}, \code{max_leaf_node_count},
\srcpath{max_lower_node_count}, \srcpath{max_upper_node_count} pre-allocate the
sparse tree (leaf $\to$ lower $\to$ upper nodes) so that the CUDA graph can
be captured once. They are \emph{totals across all envs}, not per env. Newton
requires \code{upper <= lower <= leaf <= active}. The committed values
$2^{14}/2^{13}/2^{12}/2^{10}$ were found by hand for 4 envs on 6~GB; a
``capacity exceeded'' error names the tier to raise, and a
``Failed to create volume'' plus CUDA error 700 means the caps were raised
too far for the memory available.

\paragraph{Numerics.}
\code{max_iterations=24} and \code{tolerance=1e-4} bound the per-step
rheology solve. \code{solver="auto"} lets Newton pick Gauss--Seidel for the
\code{Q1} velocity basis. \srcpath{warmstart_mode="auto"} reuses last step's
solution as a starting guess. \code{strain_basis="P0"} (piecewise constant
strain per cell), \code{velocity_basis="Q1"} (trilinear velocities) and
\srcpath{transfer_scheme="apic"} (affine particle-in-cell transfers, less
numerical damping than PIC) are the standard granular preset also used by
Isaac Lab's own \code{ur10_particle_push} task.

\paragraph{Colliders.}
\srcpath{collider_basis="pic27"} samples collider velocities with a 27-point
stencil; \srcpath{collider_velocity_mode="forward"} uses the colliders'
velocities at the start of the step. \srcpath{project_outside_colliders=False}
is mandatory on a coupled entry: it disables a particle-level push-out pass
that would fight the coupler.

\paragraph{The material.}
Particles carry an \code{MPMParticleMaterialCfg}
(\srcpath{sim/spawners/mpm/mpm_cfg.py}). Its fields: \code{density}
[kg/m$^3$], \code{young_modulus} [Pa], \code{poisson_ratio},
\code{viscosity} [Pa$\cdot$s], \code{friction} (dimensionless, $\approx
\tan\phi$ of the internal friction angle), \code{damping} [s],
\code{yield_pressure} [Pa], \code{tensile_yield_ratio}, \code{yield_stress}
[Pa] (cohesion), \code{hardening}, \code{dilatancy}. The yield surface is
\[
  \tau_{\max}(p) \;=\; \texttt{yield\_stress} \;+\; \texttt{friction}\cdot(p - p_{\min}),
\]
so shear strength grows with pressure $p$ like dry sand. The tricycle soil
sets \code{density=1800}, \code{friction=0.84} ($=\tan 40^\circ$) and
\code{yield_pressure=1e12} (effectively incompressible), leaving the very
stiff default \code{young_modulus} so the strip behaves like packed ground
rather than a sandbox.

\paragraph{The particles.}
\code{MPMGridCfg} fills an axis-aligned box with a lattice. Its resolution
per axis is $\lceil \texttt{particles\_per\_cell}\cdot \text{extent}/\texttt{voxel\_size}\rceil$,
so the $1.5\times0.6\times0.06$~m strip at 5~cm voxels gives $30\times12\times2 = 720$
particles per env, each placed at a cell centre
(\srcpath{particle_placement="cell_center"}) with mass = cell volume $\times$
density.

\begin{whybox}{the soil needs its own floor}
The MPM solver only sees colliders that belong to \emph{its} coupler entry.
The world ground plane belongs to the rigid entry, so to the soil it does not
exist and particles would fall forever. \code{tricycle_env_cfg.py} therefore
adds \code{MPMGround}, an invisible kinematic slab just below the strip, and
lists it as the MPM entry's only body.
\end{whybox}

\section{Coupling the two solvers}
\label{sec:coupling}

\begin{newtonbox}{one Newton model, two solvers}
Newton's experimental \emph{coupled solver} takes a single \code{Model} and a
list of \textbf{entries}. Each entry owns a subset of bodies, joints, shapes
and particles (a \code{ModelView}) and runs its own solver on that subset.
\textbf{Proxy mappings} expose selected bodies of one entry inside another
as \emph{virtual proxies}: massive ghost bodies whose pose is copied from the
source each step and whose accumulated contact impulses are sent back as
feedback forces. Newton calls this \emph{lagged-impulse} coupling because the
feedback is applied one step late.
\end{newtonbox}

Isaac Lab wraps this as \code{CouplerProxyCfg}
(\srcpath{source/isaaclab_contrib/isaaclab_contrib/coupling/coupler_cfg.py}):

\begin{itemize}[leftmargin=1.6em]
  \item \code{entries}: a list of \code{CouplerEntryCfg}. Each has a
        \code{name}, a \code{solver_cfg}, a list of \code{bodies} (regexes
        matched against full Newton body labels, descendants included),
        \code{all_particles}, \srcpath{include_static_shapes} (whether the entry
        owns shapes with no body, such as the ground plane),
        \srcpath{include_child_joints}, \code{substeps} (how many equal
        sub-steps the entry takes inside one coupled step) and
        \code{in_place} (step into the same state buffer; required
        \code{True} for MPM).
  \item \code{proxies}: a list of \code{CouplerProxyMappingCfg} with
        \code{source}, \code{destination}, \code{bodies},
        \code{mode} (\code{"lagged"} or \code{"staggered"}), \srcpath{mass_scale} and
        \code{collision_pipeline}.
  \item \code{iterations}: proxy relaxation passes per coupled step.
\end{itemize}

\begin{figure}[H]
\centering
\begin{tikzpicture}[
  ent/.style={draw,rounded corners=3pt,minimum width=5.6cm,minimum height=2.6cm,align=left,font=\small,fill=white},
  arr/.style={-{Stealth[length=2mm]},thick},
]
\node[ent,draw=isaacblue,fill=isaacblue!5] (rig) at (0,0) {\textbf{entry ``rover''} (MJWarp, 2 substeps)\\[2pt]
  bodies: \srcpath{/World/envs/env_.*/Tricycle}\\ + all static shapes (ground plane)\\[2pt] solves: car dynamics, wheel drives,\\ car vs ground contacts};
\node[ent,draw=newtonpurple,fill=newtonpurple!5] (mpm) at (7.6,0) {\textbf{entry ``soil''} (implicit MPM, 1 substep)\\[2pt]
  bodies: \srcpath{/World/envs/env_.*/MPMGround}\\ + all particles\\[2pt] solves: soil flow, soil vs slab,\\ soil vs \emph{proxy} wheels};
\draw[arr,isaacblue] ([yshift=5mm]rig.east) -- node[above,font=\scriptsize,align=center]{proxy: wheel poses \& velocities\\ (mass $\times$ 10)} ([yshift=5mm]mpm.west);
\draw[arr,newtonpurple] ([yshift=-5mm]mpm.west) -- node[below,font=\scriptsize,align=center]{lagged feedback: soil impulses\\ on the wheels, next step} ([yshift=-5mm]rig.east);
\end{tikzpicture}
\caption{The two coupler entries and the one proxy mapping in \code{TricycleEnvCfg.__post_init__}. Only the three wheel bodies are mirrored into the soil solver.}
\label{fig:coupler}
\end{figure}

What happens each 10~ms physics step: the rigid entry advances the car by two
5~ms sub-steps using last step's soil feedback; the wheels' new poses and
velocities are copied onto their proxies inside the MPM view; the MPM entry
advances the soil one 10~ms step, colliding particles against the proxy
wheels and the slab; the impulses the particles exerted on the proxies are
stored and applied to the real wheels during the \emph{next} rigid step.

\begin{whybox}{\code{mass_scale} is the stability knob}
A proxy body is given the source body's mass times \code{mass_scale}. A
heavy proxy barely moves when soil pushes on it, so the soil solve stays
well-posed and the feedback impulses are smooth; the real wheel then receives
those impulses one step late. If the scale is too small the lagged feedback
can oscillate (the car ``chatters'' on the strip); too large and the soil
feels the wheel as an immovable wall. The tricycle uses 10 through
\srcpath{TricycleEnvCfg.proxy_mass_scale}.
\end{whybox}

\begin{warnbox}{rules the coupler enforces}
\begin{itemize}[leftmargin=1.4em,nosep]
  \item An MPM entry must have \code{in_place=True}; the coupler raises
        \code{ValueError} otherwise (\code{coupler.py}, \code{_validate_config}).
  \item A proxy \emph{source} entry may not be \code{in_place} in lagged mode
        (Newton, \srcpath{solver_coupled_proxy.py}). The rigid entry is not, so
        this is satisfied.
  \item \srcpath{project_outside_colliders} must be \code{False} on a coupled
        MPM entry.
  \item Newton's proxy coupler supports at most two entries.
\end{itemize}
\end{warnbox}

\section{The collision pipeline and \texttt{soft\_contact\_max=0}}

With \srcpath{use_mujoco_contacts=False}, rigid contacts (car against ground
plane) come from Newton's \code{CollisionPipeline}, configured by
\code{NewtonCollisionPipelineCfg} (\srcpath{physics/newton_collision_cfg.py}).
The only field the tricycle sets is \code{soft_contact_max=0}. ``Soft
contacts'' are particle-versus-shape contacts generated by the generic
pipeline; their default allocation is \code{shape_count * particle_count},
which for thousands of particles is a huge buffer, and the MPM solver handles
particle contacts itself anyway. Setting the maximum to zero disables that
allocation.

\section{Time stepping and CUDA graphs}

\code{SimulationCfg.dt = 1/100} is the physics step; \srcpath{NewtonCfg.num_substeps=1}
means one coupled solve per step (the rigid entry's own \code{substeps=2}
happens inside it). \code{use_cuda_graph=True} records the whole step as a
CUDA graph the first time and replays it afterwards, which is where most of
the speed comes from; it also explains why buffer sizes (grid capacities,
\code{njmax}) must be fixed up front. \code{decimation=2} in the env cfg
means the policy acts every second physics step, i.e.\ at 50~Hz.

\section{Why only a handful of envs fit}

Each env carries 720 particles plus its own sparse grid; the grid tiers are
shared caps. The GPU work per step is dominated by the number of \emph{active
grid cells} times envs, and the 6~GB card runs out of memory well before the
GPU runs out of compute. The runbook's numbers (4 envs, 90--160 steps/s) are
the practical ceiling for this card; the same configuration on a large GPU
simply needs the four caps raised in proportion to \code{--num_envs}.
